\documentclass{article}

\usepackage{amsmath,amssymb}
\usepackage{xurl}
\usepackage[hidelinks]{hyperref}
\setlength{\emergencystretch}{2em}

% Shared commands for notes in docs/paper-gaps/.

\DeclareUnicodeCharacter{2080}{\ifmmode{}_{0}\else\textsubscript{0}\fi}
\DeclareUnicodeCharacter{2081}{\ifmmode{}_{1}\else\textsubscript{1}\fi}
\DeclareUnicodeCharacter{2082}{\ifmmode{}_{2}\else\textsubscript{2}\fi}
\DeclareUnicodeCharacter{2099}{\ifmmode{}_{n}\else\textsubscript{n}\fi}

\newcommand{\Lean}{\textsc{Lean}}
\ExplSyntaxOn
\NewDocumentCommand{\leanid}{m}{
  \ifmmode
    \textsf{\detokenize{#1}}
  \else
    \group_begin:
    \urlstyle{sf}
    \tl_set:Nn \l_tmpa_tl {#1}
    \tl_replace_all:Nnn \l_tmpa_tl {\_} {_}
    \exp_args:NV \path \l_tmpa_tl
    \group_end:
  \fi
}
\ExplSyntaxOff
\providecommand{\tr}{\operatorname{tr}}
\newcommand{\tnleanIssueBase}{https://github.com/LionSR/TNLean/issues/}
\newcommand{\tnleanPrBase}{https://github.com/LionSR/TNLean/pull/}
\newcommand{\tnleanDocBase}{https://sirui-lu.com/TNLean/docs/}
\newcommand{\ghissue}[1]{\href{\tnleanIssueBase#1}{GitHub issue \##1}}
\newcommand{\ghpr}[1]{\href{\tnleanPrBase#1}{PR \##1}}
\newcommand{\doclink}[1]{\href{\tnleanDocBase#1.html}{\path{#1}}}

% Dirac notation and the identity operator, as in blueprint/src/macros/common.tex.
% \providecommand keeps note-local definitions authoritative.
\usepackage{dsfont}
\providecommand{\ket}[1]{|#1\rangle}
\providecommand{\bra}[1]{\langle #1|}
\providecommand{\braket}[2]{\langle #1 | #2 \rangle}
\providecommand{\ketbra}[2]{|#1\rangle\!\langle #2|}
\providecommand{\Id}{\mathds{1}}
\providecommand{\one}{\mathds{1}}
\providecommand{\C}{\mathbb{C}}
\providecommand{\MN}[1]{M_{#1}(\C)}
\providecommand{\MD}{M_D(\C)}

% External referencing for paper sources.  The build helper writes sanitized
% auxiliary files under build/paper-aux/ and excludes duplicated source labels.
\usepackage{xr}

% Import a generated paper auxiliary file with a caller-supplied label prefix.
% The candidate paths support builds from docs/paper-gaps/, the repository root,
% and build/paper-aux/ itself.
\newcommand{\importpaperaux}[2]{%
  \IfFileExists{../../build/paper-aux/#2.aux}{%
    \externaldocument[#1]{../../build/paper-aux/#2}%
  }{%
    \IfFileExists{build/paper-aux/#2.aux}{%
      \externaldocument[#1]{build/paper-aux/#2}%
    }{%
      \IfFileExists{#2.aux}{%
        \externaldocument[#1]{#2}%
      }{}%
    }%
  }%
}

% General external reference macros
\newcommand{\paperref}[2]{\ref{#1-#2}}
\newcommand{\papereqref}[2]{(\ref{#1-#2})}

% CPSV16 source (1606.00608 / MPDO-22-12-17-2.tex) external document and shortcuts
\importpaperaux{cpsv16-}{1606.00608/MPDO-22-12-17-2-tnlean}
\newcommand{\cpsvref}[1]{\paperref{cpsv16}{#1}}
\newcommand{\cpsveqref}[1]{\papereqref{cpsv16}{#1}}

% Verdict marker: \gapnote{<kind>}{<status>}. Kind is the policy
% classification (clarification, local-correction, scope-restriction,
% unfaithful, false-source, open-gap); status is open, wip (elimination
% underway), resolved, or historical. Machine-read by the site index;
% prints a verdict line.
\newcommand{\gapnote}[2]{\par\noindent\textbf{Verdict:} #1 (#2).\par}


\title{The Vertical Direction in Proposition 4.13 Is the Index-Exchanged Tensor}
\author{}
\date{2026-07-11}

\begin{document}
\maketitle
\gapnote{unfaithful}{resolved}

\section{Source definition of the vertical direction}

Proposition~4.13, whose source label is \path{Prop:IV.12}, concerns a tensor
viewed in the vertical direction
\cite[Proposition~4.13]{Cirac2016MPDO_arXiv}. The source defines this direction
by considering the tensor ``as generating MPV by concatenating it vertically,
and thus by exchanging the notion of physical and virtual indices''
\cite[line~943]{Cirac2016MPDO_arXiv}. The local source is
\path{Papers/1606.00608/MPDO-22-12-17-2.tex}.

Let $M^{ij}_{ab}$ denote the original tensor entries, where $(i,j)$ is the
physical pair and $(a,b)$ is the horizontal bond pair. The vertically viewed
tensor $\widetilde M$ has letters indexed by $(a,b)$, and each letter acts on
the physical space with matrix entries
\begin{align}
    (\widetilde M_{ab})_{ij}
        &=M^{ij}_{ab}.
    \label{eq:cpgsv17_vertical_tensor_definition_vertical_letters}
\end{align}
Thus vertical viewing exchanges the roles of the physical and horizontal
virtual index pairs.

The conclusion of Proposition~4.13, restated at source lines~1863--1870 in the
formula labelled \path{UMU-appendix}, is
\begin{align}
    U\widetilde M U^\dagger
        &=
        \bigoplus_\alpha\mu_\alpha\otimes M_\alpha.
    \label{eq:cpgsv17_vertical_tensor_definition_vertical_canonical_form}
\end{align}
Here $U$ acts on the physical space, each $\mu_\alpha$ is a positive diagonal
matrix, and the family $\{M_\alpha\}_\alpha$ is a basis of normal tensors for
$\widetilde M$ \cite[Proposition~4.13]{Cirac2016MPDO_arXiv}. The
renormalization fixed-point characterization in Theorem~4.14 uses the
multiplicity coefficients
\begin{align}
    m_\alpha
        &=\tr(\mu_\alpha).
    \label{eq:cpgsv17_vertical_tensor_definition_multiplicity_coefficient}
\end{align}
This use occurs at source line~1956 of
Ref.~\cite{Cirac2016MPDO_arXiv}.

\section{The former diagonal-tensor surrogate}

Before the realignment, the formal vertical canonical-form predicate placed
the conclusion of Proposition~4.13 on the diagonal MPS tensor
$i\mapsto M^{ii}$. It required that diagonal tensor to generate the same
matrix product vector family as a block-diagonal tensor formed by repeating
basis blocks with positive scalar weights. Both sides used the physical
alphabet $\{0,\ldots,d-1\}$, and their matrices acted on the horizontal bond
space. This orientation is the opposite of
Eq.~\ref{eq:cpgsv17_vertical_tensor_definition_vertical_letters}, in which
letters are horizontal bond pairs and matrices act on the physical space.

The diagonal surrogate also cannot retain the multiplicity coefficients in
Eq.~\ref{eq:cpgsv17_vertical_tensor_definition_vertical_canonical_form}.
Consider physical dimension $d=2$, horizontal bond dimension $D=1$, and
\begin{align}
    M^{ij}
        &=\frac{\delta_{ij}}{2}.
    \label{eq:cpgsv17_vertical_tensor_definition_scalar_mpdo}
\end{align}
This tensor is a matrix product density operator in canonical form. Its
vertically viewed tensor has the single physical-space letter
\begin{align}
    \widetilde M_{11}
        &=\frac{1}{2}I_2.
    \label{eq:cpgsv17_vertical_tensor_definition_scalar_vertical_letter}
\end{align}
Its decomposition in
Eq.~\ref{eq:cpgsv17_vertical_tensor_definition_vertical_canonical_form} has
one block with
\begin{align}
    \mu
        &=
        \operatorname{diag}\left(\frac{1}{2},\frac{1}{2}\right),
    \label{eq:cpgsv17_vertical_tensor_definition_twofold_weight}\\
    \tr(\mu)
        &=1.
    \label{eq:cpgsv17_vertical_tensor_definition_twofold_multiplicity_sum}
\end{align}
The block occurs with multiplicity two.

By contrast, the diagonal tensor generates the scalar family
\begin{align}
    V^{(N)}
        &=2^{-N}.
    \label{eq:cpgsv17_vertical_tensor_definition_diagonal_scalar_family}
\end{align}
A flattening with two positive weights $\omega_1$ and $\omega_2$ would require
\begin{align}
    \omega_1^N+\omega_2^N
        &=2^{-N}
        \qquad\text{for every }N\geq1.
    \label{eq:cpgsv17_vertical_tensor_definition_two_weight_power_sum}
\end{align}
The cases $N=1$ and $N=2$ give
\begin{align}
    \omega_1+\omega_2
        &=\frac{1}{2},
    \label{eq:cpgsv17_vertical_tensor_definition_power_sum_first}\\
    \omega_1^2+\omega_2^2
        &=\frac{1}{4},
    \label{eq:cpgsv17_vertical_tensor_definition_power_sum_second}\\
    2\omega_1\omega_2
        &=
        (\omega_1+\omega_2)^2
        -(\omega_1^2+\omega_2^2)
        =0.
    \label{eq:cpgsv17_vertical_tensor_definition_weight_product_zero}
\end{align}
Equation~\ref{eq:cpgsv17_vertical_tensor_definition_weight_product_zero}
contradicts positivity of both weights. The surrogate can therefore represent
this tensor only with multiplicity one, so it loses the multiplicity data
$m_\alpha$ required at source line~1956. The diagonal-tensor statement is not
a corollary of Proposition~4.13, and proving it would not formalize that
proposition.

\section{The realigned formal statement}

The realignment tracked in
\href{https://github.com/LionSR/TNLean/issues/3712}{GitHub issue \#3712}
introduces the tensor from
Eq.~\ref{eq:cpgsv17_vertical_tensor_definition_vertical_letters} as a formal
object and states the vertical canonical-form predicate directly on it. The
predicate provides a basis of normal tensors $\{M_\alpha\}_\alpha$ for
$\widetilde M$, positive diagonal matrices
\begin{align}
    \mu_\alpha
        &=
        \operatorname{diag}\bigl(        \omega_{\alpha,0},
            \ldots,
            \omega_{\alpha,r_\alpha-1}
        \bigr),
    \label{eq:cpgsv17_vertical_tensor_definition_diagonal_sector_weight}\\
    r_\alpha
        &\geq1,
    \label{eq:cpgsv17_vertical_tensor_definition_positive_multiplicity}
\end{align}
and a coisometry $U$ onto the retained nonzero-sector space. These data satisfy
Eq.~\ref{eq:cpgsv17_vertical_tensor_definition_vertical_canonical_form}
letter by letter for every horizontal bond pair.

The coisometry satisfies
\begin{align}
    UU^\dagger
        &=I
    \label{eq:cpgsv17_vertical_tensor_definition_coisometry}
\end{align}
on the retained space, so $U^\dagger$ is its isometric inclusion into the
original physical space. The general canonical-form construction in
Ref.~\cite{Cirac2016MPDO_arXiv} allows $\sum_kD_k<D$ and therefore permits a
zero orthogonal complement at source lines~214--225. Consequently,
$U^\dagger U$ is generally the support projection rather than the identity.
This correction is documented in
\path{cpgsv17_vertical_isometry_zero_sector.tex}.

Proposition~4.13 does not print a separate reverse identity. Its preceding
assertion that the vertical tensor is itself in canonical form supplies exact
reconstruction and requires the omitted complement to be a zero sector
\cite[Proposition~4.13]{Cirac2016MPDO_arXiv}. Because each $\mu_\alpha$ is
diagonal,
\begin{align}
    \mu_\alpha\otimes M_\alpha
        &=
        \bigoplus_{k=0}^{r_\alpha-1}
        \omega_{\alpha,k}M_\alpha.
    \label{eq:cpgsv17_vertical_tensor_definition_weighted_sector_sum}
\end{align}
Hence the right-hand side of
Eq.~\ref{eq:cpgsv17_vertical_tensor_definition_vertical_canonical_form} is a
block-diagonal tensor indexed by pairs $(\alpha,k)$. It retains both the
multiplicities $r_\alpha$ and the weights $\omega_{\alpha,k}$. In particular,
\begin{align}
    m_\alpha
        &=
        \tr(\mu_\alpha)
        =
        \sum_{k=0}^{r_\alpha-1}\omega_{\alpha,k}.
    \label{eq:cpgsv17_vertical_tensor_definition_recovered_multiplicity}
\end{align}
The formal one-site actions used in the invariant-projection step are
identified with left and right multiplication on the letters of
$\widetilde M$.

The diagonal-tensor identities remain as auxiliary results. The matrix product
vectors of $i\mapsto M^{ii}$ are the diagonal entries
\begin{align}
    \bra{\sigma}\rho^{(N)}(M)\ket{\sigma}
        &\geq0.
    \label{eq:cpgsv17_vertical_tensor_definition_diagonal_density_entry}
\end{align}
These entries provide positivity input for the proof of Proposition~4.13, but
they no longer state its conclusion.

\section{Relation to the neighboring notes}

The note \path{cpgsv17_vertical_cf_grouping.tex} explains that the positive
diagonal coefficients in
Eq.~\ref{eq:cpgsv17_vertical_tensor_definition_vertical_canonical_form} arise
from the vertical canonical decomposition rather than from flattening
horizontal weights. The note
\path{cpgsv17_vertical_diagonal_restriction.tex} shows that diagonal
restriction of horizontal blocks does not preserve normality. Neither note
addresses the definitional mismatch corrected here: the former predicate
placed the conclusion of Proposition~4.13 on the wrong tensor.

The literal horizontal-to-vertical implication is now proved by
\path{MPOTensor.verticalCF_of_cpsvCanonicalForm} in
\path{TNLean/MPS/MPDO/CPSVVerticalCanonicalForm.lean}. It assumes the literal
CPSV canonical form and the MPDO positivity hypothesis of Proposition~4.13 in
Ref.~\cite{Cirac2016MPDO_arXiv}. The earlier theorem
\path{MPOTensor.verticalCF_of_horizontalCF} in
\path{TNLean/MPS/MPDO/VerticalCanonicalForm.lean} remains available as a
separate result under the stronger normalized BNT-refined hypothesis
\path{MPOTensor.IsHorizontalCF}.

\paragraph{Verdict.}
The former predicate had a definitional mismatch: it placed the conclusion of
Proposition~4.13 on the diagonal tensor $i\mapsto M^{ii}$ rather than on the
vertically viewed tensor $\widetilde M$ defined at source line~943 of
Ref.~\cite{Cirac2016MPDO_arXiv}. The example with $d=2$ and $D=1$ shows that
the diagonal surrogate cannot retain the coefficients
$m_\alpha=\tr(\mu_\alpha)$ from
Eq.~\ref{eq:cpgsv17_vertical_tensor_definition_vertical_canonical_form}.
A proof of that surrogate would therefore not have formalized the proposition.

The realignment restores the source-faithful vertical tensor definition. The
later correction in
\path{cpgsv17_vertical_isometry_zero_sector.tex} expresses the source's
one-sided condition as a coisometry onto the retained nonzero sectors and
retains exact reconstruction. The horizontal-to-vertical implication is now
proved with this corrected definition under the literal CPSV canonical-form
and MPDO hypotheses, while the stronger normalized BNT-refined theorem remains
available separately. The phase-class grouping and positive diagonal
coefficients used in the proof are detailed in
\path{cpgsv17_vertical_cf_grouping.tex}.

This work completes Proposition~4.13. It does not prove the positive-fusion
statement in Appendix~C.4 or the renormalization fixed-point characterization
in Theorem~4.14 of Ref.~\cite{Cirac2016MPDO_arXiv}.

\bibliographystyle{alpha}
\bibliography{references}

\end{document}
